\section*{Capítulo \ref{ch:b2:angiosperm-reproduction} --- Reprodução sexuada das plantas com flores}

\begin{solution}{exo:b2:angiosperm-reproduction:1}
Sépalas (cálice), pétalas (corola), \omterm{def:b2:angiosperm-reproduction:flower}{estames} (androceu), \omterm{def:b2:angiosperm-reproduction:flower}{carpelos}
(gineceu). \omterm{def:b2:angiosperm-reproduction:flower}{Estame}: filete e antera. \omterm{def:b2:angiosperm-reproduction:flower}{Carpelo}: \omterm{def:b2:angiosperm-reproduction:flower}{estigma}, estilete, \omterm{def:b2:angiosperm-reproduction:flower}{ovário}
com os \omterm{def:b2:plant-life-cycles:heterospory}{óvulos}. Todas essas partes são \omterm{def:b2:plant-life-cycles:cycles}{esporófito} \omterm{def:b2:meiosis-heredity:meiosis}{diploide}; apenas o
\omterm{prop:b2:angiosperm-reproduction:pollen}{grão de pólen} (dentro da antera) e o \omterm{prop:b2:angiosperm-reproduction:embryosac}{saco embrionário} (dentro do
\omterm{def:b2:plant-life-cycles:heterospory}{óvulo}) são \omterm{def:b2:plant-life-cycles:cycles}{gametófito} \omterm{def:b2:meiosis-heredity:meiosis}{haploide}.
\end{solution}

\begin{solution}{exo:b2:angiosperm-reproduction:2}
\omterm{def:b2:plant-life-cycles:heterospory}{Micrósporo} $n$; célula vegetativa $n$; célula espermática $n$; oosfera
$n$; núcleo polar $n$; zigoto $2n$; endosperma $3n$; \omterm{def:b2:angiosperm-reproduction:seed}{tegumento da
semente} $2n$ (tegumentos maternos); parede do \omterm{def:b2:angiosperm-reproduction:seed}{fruto} $2n$ (\omterm{def:b2:angiosperm-reproduction:flower}{ovário}
materno).
\end{solution}

\begin{solution}{exo:b2:angiosperm-reproduction:3}
Vento (gramíneas, carvalho, bétula, aveleira, tanchagem): \omterm{def:b2:angiosperm-reproduction:flower}{flores}
pequenas e verdes, sem pétalas, perfume ou \omterm{def:b2:angiosperm-reproduction:pollination}{néctar}, anteras pendentes em
filetes longos, \omterm{def:b2:angiosperm-reproduction:flower}{estigmas} plumosos, enormes quantidades de \omterm{def:b2:plant-life-cycles:heterospory}{pólen} liso e
seco, floração antes das folhas. Abelha (trevo, sálvia, dedaleira,
lavanda, macieira): pétalas coloridas, muitas vezes com guias
ultravioleta, plataforma de pouso, perfume, \omterm{def:b2:angiosperm-reproduction:pollination}{néctar}, \omterm{def:b2:plant-life-cycles:heterospory}{pólen} pegajoso e
esculpido em quantidade moderada, floração durante o dia.
\end{solution}

\begin{solution}{exo:b2:angiosperm-reproduction:4}
As duas células espermáticas de um mesmo \omterm{thm:b2:angiosperm-reproduction:tube}{tubo polínico} se fundem, uma
com a oosfera (zigoto, $2n$, o embrião) e a outra com a célula central
(núcleo do endosperma, $3n$, a reserva de alimento). O endosperma só
começa a crescer depois da fecundação, de modo que a planta abastece
apenas os \omterm{def:b2:plant-life-cycles:heterospory}{óvulos} que contêm um embrião, ao contrário da gimnosperma,
que constrói sua reserva antes.
\end{solution}

\begin{solution}{exo:b2:angiosperm-reproduction:5}
$25/1.2 = \qty{21}{h}$. Os cabelos que emergem nos dias 5, 6 e 7
encontram \omterm{def:b2:plant-life-cycles:heterospory}{pólen}; os dos dias 8 a 12, não: $3/8$ dos cabelos são
polinizados, e a espiga fica cinco oitavos estéril.
\end{solution}

\begin{solution}{exo:b2:angiosperm-reproduction:6}
$S_1S_2\times S_1S_2$: 0. $S_1S_2\times S_1S_3$: $\tfrac12$ (só o \omterm{def:b2:plant-life-cycles:heterospory}{pólen}
$S_3$ cresce). $S_1S_2\times S_3S_4$: todo. Descendentes do segundo
cruzamento: oosferas $S_1$ ou $S_2$, células espermáticas $S_3$:
$S_1S_3$ e $S_2S_3$, metade de cada.
\end{solution}

\begin{solution}{exo:b2:angiosperm-reproduction:7}
Uma planta tem 2 dos $n$ \omterm{def:b2:meiosis-heredity:vocabulary}{alelos} e rejeita o \omterm{def:b2:plant-life-cycles:heterospory}{pólen} que leve qualquer um
deles: fração compatível $(n - 2)/n$. Um \omterm{def:b2:meiosis-heredity:vocabulary}{alelo} novo não é rejeitado por
nenhum pistilo, de modo que seu \omterm{def:b2:plant-life-cycles:heterospory}{pólen} gera mais \omterm{def:b2:angiosperm-reproduction:seed}{sementes} que a média e
ele se espalha até ficar tão comum quanto os outros; a mesma vantagem
protege todo \omterm{def:b2:meiosis-heredity:vocabulary}{alelo} raro da perda. A seleção acumula, portanto, \omterm{def:b2:meiosis-heredity:vocabulary}{alelos},
e as populações naturais de espécies autoincompatíveis têm dezenas
deles.
\end{solution}

\begin{solution}{exo:b2:angiosperm-reproduction:8}
Os dois \omterm{prop:b2:angiosperm-reproduction:embryosac}{núcleos polares} são cópias de um mesmo \omterm{def:b2:plant-life-cycles:heterospory}{megásporo}, de modo que a
célula central é $YY$ ou $yy$ (metade de cada); a célula espermática é
$Y$ ou $y$ (metade de cada): endosperma $YYY$, $YYy$, $Yyy$, $yyy$, a
$\tfrac14$ cada; três quartos amarelos. Polinizada por $yy$: $YYy$ e
$yyy$, metade de cada --- metade dos grãos amarelos.
\end{solution}

\begin{solution}{exo:b2:angiosperm-reproduction:9}
Grãos de cevada são cortados ao meio; as metades com embrião digerem
seu amido, as metades sem embrião não, a menos que se acrescente
giberelina, caso em que o digerem. (a) O embrião é a fonte do sinal, e
não da enzima. (b) A giberelina é o sinal, suficiente por si só. (c) A
amilase é produzida pela camada de aleurona do endosperma, que responde
ao hormônio. A metade sem embrião separa o sinal da resposta: sem ela,
não se poderia saber se o próprio embrião digere o amido.
\end{solution}

\begin{solution}{exo:b2:angiosperm-reproduction:10}
Vento: $10^{6}\times\qty{1}{\micro g} = \qty{1}{g}$ por \omterm{def:b2:plant-life-cycles:heterospory}{óvulo}. Animal:
\qty{1}{mg} de \omterm{def:b2:plant-life-cycles:heterospory}{pólen} mais \qty{0.5}{mg} de açúcar $= \qty{1.5}{mg}$
por \omterm{def:b2:plant-life-cycles:heterospory}{óvulo} --- cerca de setecentas vezes mais barato. A \omterm{def:b2:angiosperm-reproduction:pollination}{polinização} pelo
vento não precisa de parceiro: funciona no início da primavera, antes
que os insetos voem, em lugares frios, ventosos e abertos, à noite e
na chuva, em populações densas de uma só espécie, e não pode ser
enganada por ladrões de \omterm{def:b2:angiosperm-reproduction:pollination}{néctar} nem abandonada por um polinizador que
se extinga.
\end{solution}

\begin{solution}{exo:b2:angiosperm-reproduction:11}
As \omterm{def:b2:angiosperm-reproduction:seed}{sementes} que caem ao pé da planta-mãe pousam à sua sombra, entre
suas raízes e onde se concentram seus predadores de \omterm{def:b2:angiosperm-reproduction:seed}{sementes} e seus
patógenos específicos; quase todas morrem. Uma ave leva uma \omterm{def:b2:angiosperm-reproduction:seed}{semente} a
quilômetros (\qty{12}{km} em vinte minutos) e a deixa cair, com adubo,
numa sebe ou numa clareira; mesmo que só uma \omterm{def:b2:angiosperm-reproduction:seed}{semente} em cem caia num
bom lugar, isso é mais do que todas as \omterm{def:b2:angiosperm-reproduction:seed}{sementes} debaixo da árvore, e a
população se espalha à velocidade das aves, e não à dos \omterm{def:b2:angiosperm-reproduction:seed}{frutos} que
caem. A polpa é o preço da passagem.
\end{solution}

\begin{solution}{exo:b2:angiosperm-reproduction:12}
O encerramento obriga o \omterm{def:b2:plant-life-cycles:heterospory}{pólen} a crescer através de tecido materno, o
que permite ao pistilo testá-lo (autoincompatibilidade, rejeição de
espécies estranhas), deixar muitos tubos competirem para que o mais
rápido gere a \omterm{def:b2:angiosperm-reproduction:seed}{semente}, proteger os \omterm{def:b2:plant-life-cycles:heterospory}{óvulos} da dessecação e dos
herbívoros, e transformar o \omterm{def:b2:angiosperm-reproduction:flower}{ovário} num \omterm{def:b2:angiosperm-reproduction:seed}{fruto} que recruta animais para
a dispersão; a \omterm{prop:b2:angiosperm-reproduction:double}{dupla fecundação} passa então a abastecer apenas os
\omterm{def:b2:plant-life-cycles:heterospory}{óvulos} fecundados. Os custos: um \omterm{prop:b2:angiosperm-reproduction:pollen}{grão de pólen} precisa alcançar um
\omterm{def:b2:angiosperm-reproduction:flower}{estigma}, e não o próprio \omterm{def:b2:plant-life-cycles:heterospory}{óvulo}, de modo que são necessários um
estilete, um tubo e sua orientação, e o \omterm{def:b2:angiosperm-reproduction:seed}{fruto} é um investimento
pesado. O saldo favoreceu tanto as angiospermas que elas passaram de
nada a nove décimos das espécies vegetais em cem milhões de anos.
\end{solution}

\begin{solution}{pb:b2:angiosperm-reproduction:1}
\textbf{1.} Sobre A ($S_1S_2$): \omterm{def:b2:plant-life-cycles:heterospory}{pólen} de A $f = 0$; de B ($S_1$, $S_3$)
$f =
\tfrac12$; de C ($S_3$, $S_4$) $f = 1$.
\textbf{2.} Sobre B ($S_1S_3$): A $\tfrac12$, B 0, C $\tfrac12$. Sobre C
($S_3S_4$): A 1, B $\tfrac12$, C 0.
\textbf{3.} 50 grãos compatíveis; $0.1\times 50 = 5$ \omterm{def:b2:plant-life-cycles:heterospory}{óvulos}.
\textbf{4.} De C: 100 compatíveis, $0.1\times 100 = 10$, isto é, os
dez \omterm{def:b2:plant-life-cycles:heterospory}{óvulos}; de A: nenhum.
\textbf{5.} Os \omterm{def:b2:angiosperm-reproduction:seed}{frutos} de visitas vindas de C (10 \omterm{def:b2:angiosperm-reproduction:seed}{sementes}) são
mantidos; os de B (5 \omterm{def:b2:angiosperm-reproduction:seed}{sementes}), mantidos por pouco; os de A,
descartados. Um pomar só de A não produz nada.
\textbf{6.} Metade das \omterm{def:b2:angiosperm-reproduction:flower}{flores} frutifica (as visitadas a partir de C):
2500 \omterm{def:b2:angiosperm-reproduction:seed}{frutos} por árvore, dez vezes os 250 necessários --- uma safra
completa, e a árvore descartará o excedente.
\textbf{7.} Ela libera \omterm{def:b2:plant-life-cycles:heterospory}{pólen} compatível ao longo de toda a estação de
floração de todas as variedades; mas, se compartilhar os dois \omterm{def:b2:meiosis-heredity:vocabulary}{alelos}
$S$ com uma variedade, não poliniza nada nela.
\textbf{8.} $8\times 10^{7}/(7\times 86400) = 132$ grãos por metro
quadrado por segundo.
\textbf{9.} $132/0.25 = 530$ grãos por metro cúbico.
\textbf{10.} $132\times 10^{-4} = 0.013$ por segundo: um grão a cada
\qty{76}{s}.
\textbf{11.} $0.013\times 86400 \approx 1100$ grãos por dia; o primeiro
tubo a alcançar o \omterm{def:b2:plant-life-cycles:heterospory}{óvulo} o fecunda, e o \omterm{def:b2:plant-life-cycles:heterospory}{óvulo} fica então fechado para
os demais.
\textbf{12.} \qty{20}{h} após a \omterm{def:b2:angiosperm-reproduction:pollination}{polinização}.
\textbf{13.} \omterm{def:b2:plant-life-cycles:heterospory}{Pólen} nos dias 0--7, cabelos a partir do dia 5: só os dias
5, 6 e 7 se sobrepõem --- $3/8$ dos cabelos, e os cabelos
que emergem depois do dia 7 não recebem nada; a espiga fica em grande
parte estéril. A seca na emissão dos cabelos é a causa clássica de uma
safra de milho frustrada.
\textbf{14.} Embrião $Yy$; endosperma $YYy$.
\textbf{15.} Amarelo, pois $Y$ é \omterm{def:b2:meiosis-heredity:vocabulary}{dominante}: o caráter do genitor que
forneceu o \omterm{def:b2:plant-life-cycles:heterospory}{pólen} aparece na \omterm{def:b2:angiosperm-reproduction:seed}{semente}, sobre a planta-mãe --- é a xênia.
\textbf{16.} Célula central $YY$ ou $yy$: endosperma $YYy$ (metade) e
$yyy$ (metade).
\textbf{17.} $YYy$: 20 unidades; $Yyy$: 10 unidades; um $yyY$ com o
$Y$ vindo do pai é o mesmo \omterm{def:b2:meiosis-heredity:vocabulary}{genótipo} $Yyy$, 10 unidades. Duas doses
paternas não podem existir: uma única célula espermática fecunda a
célula central, e o endosperma tem sempre dois genomas maternos para
um paterno.
\textbf{18.} É a reserva de amido; $0.3/0.33 = \qty{91}{\%}$ do grão.
\textbf{19.} A reserva só é construída após a fecundação, de modo que
nenhum recurso vai para \omterm{def:b2:plant-life-cycles:heterospory}{óvulos} sem embrião; o \omterm{def:b2:plant-life-cycles:cycles}{gametófito} da
gimnosperma é construído antes e desperdiçado se a \omterm{def:b2:angiosperm-reproduction:pollination}{polinização}
falhar.
\textbf{20.} $0.9\times 600 = 540$ grãos por espiga e por planta;
$4320$ por metro quadrado; $4320\times 0.3 = \qty{1.3}{kg}$ de
endosperma por metro quadrado.
\textbf{21.} \qty{13}{t/ha}.
\textbf{22.} Carbono do grão $0.5\times 1.2 = \qty{0.6}{kg/m^2}$;
carbono do endosperma $0.45\times 1.3 = \qty{0.58}{kg/m^2}$: coerente
(o pequeno resto é o embrião e o envoltório).
\textbf{23.} $8\times 10^{7}/4320 \approx \num{18500}$ \omterm{prop:b2:angiosperm-reproduction:pollen}{grãos de pólen}
por grão de milho.
\textbf{24.} Sua própria nuvem de \omterm{def:b2:plant-life-cycles:heterospory}{pólen} é rala e levada pelo vento, e
seu pendão libera a maior parte do \omterm{def:b2:plant-life-cycles:heterospory}{pólen} antes que seus próprios
cabelos emerjam, de modo que muitos cabelos nunca são polinizados e a
espiga fica falhada.
\textbf{25.} Bloco A: metade das \omterm{def:b2:angiosperm-reproduction:flower}{flores} frutifica, 2500 \omterm{def:b2:angiosperm-reproduction:seed}{frutos} por
árvore, uma safra completa. Lavoura: 4320 grãos por metro quadrado,
\qty{13}{t/ha}.
\end{solution}
